\section{标准烛光截面
\label{sec:StandardCandles}
}

对撞机上单举强子产生总截面的测量为标准模型提供了基石性检验。这些相对简单的可观测量既能被高精度测量，也能在 NNLO QCD 理论中以很小的不确定度预言。
Sec.~\ref{sec:ellipse} 中，我们基于 CT14、\CTHERAII{} 与 CT18(A/X/Z)
NNLO PDF，收集了 LHC 在质心能量 $\sqrt{s} = 7$、8、13 与 14 TeV 下 $W$、$Z$ 玻色子、顶夸克对以及希格斯玻色子（经胶子-胶子融合）单举产生截面的 NNLO 理论预言。
这些理论预言取代了用上一代 CT10/CT14 PDF 所作的类似比较~\cite{Gao:2013xoa,Dulat:2015mca}，可与相应的实验测量对照。此外，我们还在 Sec.~\ref{sec:Res} 给出基于固定阶与重求和计算的 LHCb 矢量玻色子产生理论预言；在 Sec.~\ref{sec:Wcharm} 探讨 ATLAS $W\!+\!c$ 产生的预言、在 Sec.~\ref{sec:tt13} 探讨 CMS 13 TeV $t\bar{t}$ 产生的预言；并在 Sec.~\ref{sec:SeaQuest} 展示高 $x$ 固定靶 Drell-Yan 截面的预言，以迎接费米实验室 SeaQuest 实验~\cite{Aidala:2017ofy} 即将发表的结果。


\subsection{LHC 的单举总截面}
\label{sec:ellipse}

\begin{figure}[p]
	\begin{center}
  \includegraphics[width=0.44\textwidth]{images/cor_tel_ct18s_H-Z__7TeV_shifted_KP_ect.pdf}
  \includegraphics[width=0.44\textwidth]{images/cor_tel_ct18s_H-Z__8TeV_shifted_KP_ect.pdf} \\
  \includegraphics[width=0.44\textwidth]{images/cor_tel_ct18s_H-Z_13TeV_shifted_KP_ect.pdf}
  \includegraphics[width=0.44\textwidth]{images/cor_tel_ct18s_H-Z_14TeV_shifted_KP_ect.pdf} \\
  \includegraphics[width=0.44\textwidth]{images/cor_tel_ct18s_H-T__7TeV.pdf}
  \includegraphics[width=0.44\textwidth]{images/cor_tel_ct18s_H-T__8TeV.pdf} \\
  \includegraphics[width=0.44\textwidth]{images/cor_tel_ct18s_H-T_13TeV.pdf}
  \includegraphics[width=0.44\textwidth]{images/cor_tel_ct18s_H-T_14TeV.pdf} \\
	\end{center}
    \vspace{-2ex}
	\caption{LHC 7、8、13 与 14 TeV 处用 CT18 NNLO PDF 族计算的 $ggH$、$t \bar t$ 与 $Z^0$ 单举截面的 90\% C.L. 误差椭圆。
		\label{fig:corr_ellipse1}}
\end{figure}


\begin{figure}[p]
	\begin{center}
  \includegraphics[width=0.44\textwidth]{images/cor_tel_ct18s_Wp-Wm__7TeV_shifted.pdf}
  \includegraphics[width=0.44\textwidth]{images/cor_tel_ct18s_Wp-Wm__8TeV_shifted.pdf} \\
  \includegraphics[width=0.44\textwidth]{images/cor_tel_ct18s_Wp-Wm_13TeV_shifted.pdf}
  \includegraphics[width=0.44\textwidth]{images/cor_tel_ct18s_Wp-Wm_14TeV_shifted.pdf} \\
  \includegraphics[width=0.44\textwidth]{images/cor_tel_ct18s_W-Z__7TeV_shifted_KP_ect.pdf}
  \includegraphics[width=0.44\textwidth]{images/cor_tel_ct18s_W-Z__8TeV_shifted_KP_ect.pdf} \\
  \includegraphics[width=0.44\textwidth]{images/cor_tel_ct18s_W-Z_13TeV_shifted_KP_ect.pdf}
  \includegraphics[width=0.44\textwidth]{images/cor_tel_ct18s_W-Z_14TeV_shifted_KP_ect.pdf} \\
	\end{center}
	\caption{同图~\ref{fig:corr_ellipse1}，但为
$W^+$、$W^-$ 与 $Z^0$ 单举截面。
		\label{fig:corr_ellipse2}}
\end{figure}


在本工作中，顶夸克与希格斯玻色子的质量分别取
$m_t^\mathit{pole}=173.3$ GeV 与 $m_H=125$ GeV。
$W$ 与 $Z$ 单举截面
（乘以衰变到某一带电轻子味的分支比）
用 \texttt{Vrap} v0.9 程序~\cite{Anastasiou:2003ds,Anastasiou:2003yy} 按 QCD NNLO 计算，重正化与因子化标度（$\mu_R$ 与 $\mu_F$）取矢量玻色子不变质量。
单举顶夸克对总截面借助程序 \texttt{Top++}
v2.0~\cite{Czakon:2013goa,Top++} 以 NNLO+NNLL 精度计算，QCD 标度按 \texttt{Top++} 框架的默认设为顶夸克质量~\cite{Czakon:2017wor}。
经由胶子-胶子融合的希格斯玻色子截面用 \texttt{iHixs} v1.3
程序~\cite{Anastasiou:2011pi} 按 QCD NNLO 计算，
采用含有限顶夸克质量修正的重夸克有效理论（HQET），QCD 标度取希格斯玻色子质量。

\begin{figure}[tb]
	\begin{center}
		\includegraphics[width=0.65\textwidth]{images/cor_tel_ct18s_W-Z__7TeV_shifted_IC_TH.pdf}
	\end{center}
	\vspace{-2ex}
	\caption{$\sqrt{s}=7$ TeV 处与 ATLAS 7 TeV $W/Z$ 数据（Exp.~ID=248）相关的 $W$、$Z$ 玻色子总产生截面理论预言。这里还纳入了若干包含内禀粲（IC）PDF 的计算——或基于 BHPS 价型模型（三种归一化）或基于海型模型（两种归一化）——叠加在 CT14 上，见文献~\cite{Hou:2017khm} 的描述。
		\label{fig:IC}}
\end{figure}

图~\ref{fig:corr_ellipse1} 显示，LHC 经胶子-胶子融合的希格斯玻色子截面（$ggH$）与顶夸克对（$t \bar t$）截面没有显著关联，因为这两个过程由胶子 PDF 主导但处于略为不同的 $x$ 区域。$ggH$ 与 $Z$ 玻色子截面间的反关联程度随 LHC 能量升高而减弱。
另一方面，图~\ref{fig:corr_ellipse2} 显示 LHC 电弱规范玻色子截面彼此高度关联。
一般而言，CT18 的预言更接近 \CTHERAII，最大差异出现在 CT18Z 与 CT14 之间。
此外，对电弱规范玻色子产生，CT18X 预言更接近 CT18Z（参见图~\ref{fig:corr_ellipse2}）；但对 $ggH$ 或 $t \bar t$ 单举截面并非如此（参见图~\ref{fig:corr_ellipse1}）。

图~\ref{fig:corr_ellipse2} 下半部分 $W$、$Z$ 截面误差椭圆的相对位置，可以联系 Sec.~\ref{sec:CT18ZvsCT18} 讨论的 CT18、A、X、Z 组合之间奇异性及其他 PDF 的差异来理解。总体上所有 $W$-$Z$ 椭圆取向相似：沿半短轴的方向——对应 $W$ 与
$Z$ 产生截面的相对差——最直接地联系着奇异 PDF。
$s$-PDF 与 $W^\pm/Z$ 截面比的关联最早在 CTEQ6.6 分析中被指出~\cite{Nadolsky:2008zw}。
基于
CT18A 与 CT18Z 的理论预言都同样朝该方向移动。与此同时，CT18X 尤其是 CT18Z 沿半长轴［``$\sigma(Z)\!+\!\sigma(W)$ 方向''］明显偏移，这与 $x<10^{-2}$ 的胶子更相关，文献~\cite{Nadolsky:2008zw} 同样指出过这一点。
CT18Z 与 A 在垂直方向上的紧密排列，与这两种拟合所得奇异分布的相似性密切相关。

考察非微扰粲的纳入是否会显著改变这些理论预言是值得的，尤其对电弱玻色子产生。文献~\cite{Ball:2017nwa} 提出，HERA 合并数据（Exp.~ID=160）与 ATLAS 7 TeV $W/Z$ 数据之间的张力要求初始标度 $Q_0$ 处的粲 PDF 被独立参数化。这种{\it 符号不定}、形状任意的非微扰粲由此用 NNPDF 合作组独特的神经网络方法实现为质子结构中区别于微扰粲的``拟合粲''贡献。内禀粲问题——包括其在微扰 QCD 中的动力学起源——也被 CT 研究，最近见于文献~\cite{Hou:2017khm}：它把正定的非微扰粲作为显式扭度为 2 的内禀 PDF，参照多种模型，在标度 $Q_0=m_c$ 处作为粲微扰演化的边界条件。

循此工作，图~\ref{fig:IC} 给出 7 TeV 总 $W,Z$ 产生截面的理论预言，类比图~\ref{fig:corr_ellipse2} 左起第三面板，但包含若干 IC 情景。CTEQ6.6 分析的关联研究~\cite{Nadolsky:2008zw} 表明：中心点 $\{\sigma_W, \sigma_Z\}$ 主要通过增大相关区域 $x\approx 0.01$ 的胶子 PDF 而朝图~\ref{fig:corr_ellipse2} 中的方向 $G$ 移动；主要通过在相近 $x$ 区域增大 $s$ PDF 而朝方向 $S$ 移动。

依此经验法则，CT18A 椭圆相对 CT18 的上移与纳入 ATLAS 7 TeV $W/Z$ 数据后 CT18A 奇异性增大相一致。在 CT14 中纳入 IC 使中心 CT14 的理论预言逆方向 $S$ 移动，这与 CT14 中 DIS 实验在纳入{\it 正定} IC PDF 后偏好压制奇异性相一致。CT14 IC 预言还沿方向 $G$ 移动，反映 CT14 IC 模型的胶子 PDF 大小与名义 CT14 不同。

于是，图中 IC 预言相对 CT14 等纯外赋粲预言的下移，看来是在初始标度 $Q_0$ 假定{\it 非负粲}的一般结果——如~\cite{Hou:2017khm} 所论，这会自然来自扭度-4 贡献。
原因还是这类 IC 模型对奇异 PDF 的某种压制，只会加剧与 ATLAS 7 TeV $W/Z$ 数据的张力。我们的结论是：纳入非负内禀 $c_\mathrm{IC}(x,Q=m_c)$ PDF 的标准做法不太可能化解 ATLAS 7 TeV $W/Z$ 数据与 HERA 之间的张力。内禀/拟合粲的定义及其动力学起源所涉的微妙之处相当复杂，需要更多后续分析来厘清并理解其唯象含义。


\subsection{LHC 矢量玻色子微分截面}
\label{sec:Res}

\begin{figure}[htb]
	\includegraphics[width=0.49\textwidth]{images/data_245_CT18__1_ect.pdf}
	\includegraphics[width=0.49\textwidth]{images/data_245_CT18__1-TcD_ect.pdf}
	\includegraphics[width=0.49\textwidth]{images/data_245_CT18__2_ect.pdf}
	\includegraphics[width=0.49\textwidth]{images/data_245_CT18__2-TcD_ect.pdf}
	\includegraphics[width=0.49\textwidth]{images/data_245_CT18__3_ect.pdf}
	\includegraphics[width=0.49\textwidth]{images/data_245_CT18__3-TcD_ect.pdf}
	\caption{
LHCb 7 TeV $W,Z$ 数据与 CT18 预言的比较，分别采用 NNLO（记作 CT18）、ResBos（记作 CT18 ResBos）或 ResBos2（记作 CT18 ResBos2）计算。同时给出 CT18Z NNLO 的预言。
		\label{fig:245_res}}
\end{figure}

如前所述，NNLO 计算以往已用于预言 Tevatron 与 LHC 的矢量玻色子产生数据。
过去我们也曾把此类精密数据与含多胶子发射效应的 ResBos 预言比较~\cite{Balazs:1997xd}，用于生成 CTEQ6.6、CT10 与 CT14 PDF。
因此，把矢量玻色子微分截面测量与基于 CT18（及 CT18Z）PDF 的 ResBos 和 NNLO 计算预言相比较十分重要。作为示例，我们在图~\ref{fig:245_res} 把 ResBos 预言与 LHCb 7 TeV $W,Z$ 玻色子微分分布~\cite{Aaij:2015gna} 比较。

为完整起见，同图中我们还纳入了 ResBos2 的预言——它是 ResBos 项目的更新版本，包含完整的 NNLO 修正，即本计算包含了双轻子对 Drell-Yan 产生的完整 $\alpha_s^2$ 贡献~\cite{resbos2-paper}。
相比之下，ResBos 预言只含部分 NNLO 贡献：更具体地说，重求和计算中只含 Wilson 系数 $C^{(1)}$ 而不含 $C^{(2)}$，参见文献~\cite{Balazs:1997xd}。
如图~\ref{fig:245_res} 所示，除极大快度区外，ResBos 与 ResBos2 对 LHCb 运动学的预言符合良好。
重求和与（NNLO）固定阶预言的差异来自多重软胶子辐射，其效应在大快度区趋于增长，在那里它与 LHCb 7 TeV 实验误差大小相当。关于重求和与固定阶计算差异的进一步详细讨论将另行发表。
为了解不同 PDF 会如何修改理论与 LHCb 7 TeV 数据的这些比较，同一图中我们还给出基于 CT18Z PDF 的预言，其胶子与海夸克分布不同于 CT18。

\subsection{LHC 的 $W$ 伴随粲喷注产生}
\label{sec:Wcharm}

\begin{figure}[tb]
	\begin{center}
		\fbox{\parbox[c][3.2cm][c]{0.86\linewidth}{\centering [figure unavailable -- see original paper]\\ \texttt{hepdata\_mcfm\_grid-WpCbar-ATLAS7\_68CL}}}
		\fbox{\parbox[c][3.2cm][c]{0.86\linewidth}{\centering [figure unavailable -- see original paper]\\ \texttt{hepdata\_mcfm\_grid-WmCjet-ATLAS7\_68CL}}}
	\end{center}
	\vspace{-2ex}
	\caption{ATLAS 7 TeV $W^+ + \bar{c}$-jet（左）
		与 $W^-\! +\! c$-jet（右）产生（电子与缪子道合并）的 CT18(Z) 与 \CTHERAII~NNLO 预言比较。
	CT18 PDF 不确定度按 68\% CL 评估。
		标度选择为 $\mu_R=\mu_F=M_W$。
		\label{fig:atl7wpcbar}}
\end{figure}

\begin{figure}[tb]
	\begin{center}
		\fbox{\parbox[c][3.2cm][c]{0.86\linewidth}{\centering [figure unavailable -- see original paper]\\ \texttt{CMS7Wc\_pTl25GeV\_68CL.pdf}}}
		\fbox{\parbox[c][3.2cm][c]{0.86\linewidth}{\centering [figure unavailable -- see original paper]\\ \texttt{CMS7Wc\_pTl35GeV\_68CL.pdf}}}
	\end{center}
	\vspace{-2ex}
	\caption{CT18(Z) 与 CT14$_{\textrm{HERAII}}$ 预言同 CMS 7 TeV $W+c$ 数据的比较：左为缪子道、轻子横动量割断 $p_T^l>25$ GeV；右为电子与缪子合并道、$p_T^l>35$ GeV。
		标度选择为 $\mu_R=\mu_F=M_W$。
		\label{fig:cms7wpcbar}}
\end{figure}

CT18 的 $s$ 夸克 PDF 与 CT14 在 $x < 10^{-1}$ 处的差异，主要源于新 LHCb 与 ATLAS 7 TeV 矢量玻色子产生数据的纳入。对
奇异夸克的独立约束来自那些由 $s$ 夸克起始过程有显著贡献的截面，例如 $W$ 伴随粲喷注产生~\cite{Yalkun:2019gah}。由于该过程尚未算至 NNLO，相关数据样本未纳入 CT18(Z) NNLO PDF 拟合，但把 NLO 预言（配 NNLO PDF）与数据比较仍有启发意义。

图~\ref{fig:atl7wpcbar} 分别把 \CTHERAII、CT18 与 CT18Z 预言同 ATLAS 7 TeV $W^+ + \bar{c}$-jet 与 $W^-\! +\! c$-jet 数据比较。PDF 不确定度按 68\% CL 评估并以黄色带表示。理论计算使用 \texttt{MCFM} 生成的 \texttt{APPLgrid} 表完成，并与 \texttt{MadGraph\_aMC@NLO}+\texttt{aMCfast} 交叉核对。
该计算的标度选择为 $\mu_R=\mu_F=M_W$，$\alpha_s$ 的跑动取 LHAPDF 表随 \texttt{APPLgrid} 提供的 NNLO 版本。
对共 $N_{pt}=22$ 个数据点，\CTHERAII、CT18 与 CT18Z PDF 的 $\chi^2/N_{pt}$ 分别为 0.59、0.52 与 0.41。
我们看到图~\ref{fig:atl7wpcbar} 两面板中 CT18Z 相对 CT18 的预言都向上移动。

有趣的是，对 $W^+\bar{c}$ 产生，
\CTHERAII{} 预言在大快度区甚至比 CT18(Z) 更大；而对图~\ref{fig:atl7wpcbar} 右面板的 $W^-c$ 产生，\CTHERAII{} 预言在全绘图范围内远低于 CT18(Z) 预言。大快度 $W^+c$ 截面的这一细致行为不仅反映 CT18Z 中 $s$ 夸克 PDF 的增加，也反映大 $x$ 处 $g$、$d$ 等 PDF 的某些补偿变化——后者已通过用 \texttt{ePump} 程序以 $W+c$ 截面更新 \CTHERAII\, PDF 独立核实。为完整起见，图~\ref{fig:cms7wpcbar} 也给出与 CMS 7 TeV $W+c$ 数据的类似比较。

\begin{figure}[tb]
\begin{center}
\fbox{\parbox[c][3.2cm][c]{0.86\linewidth}{\centering [figure unavailable -- see original paper]\\ \texttt{CMS13-pT-top-CT18NNLO-mt172p5\_ect.pdf}}}
\fbox{\parbox[c][3.2cm][c]{0.86\linewidth}{\centering [figure unavailable -- see original paper]\\ \texttt{CMS13-pt-top-nnlo-ct18-zaxnnlo-mt172p5\_ect.pdf}}}
\caption{左：顶夸克横向动量 $p_{T,t}$ 分布。
右：含 CT18Z、CT18X、CT18A NNLO 的未平移数据-理论图。右图中数据与理论预言都按 CT18NNLO 理论归一。
误差棒表示统计误差与总系统误差的平方和。}
\label{pTt}
\end{center}
\end{figure}

\subsection{LHC 13 TeV 顶夸克对微分分布}
\label{sec:tt13}

Sec.~\ref{sec:QualityTopData} 已展示 8 TeV ATLAS 与 CMS 微分顶产生数据（即纳入 CT18(Z) 拟合的数据）的数据-理论比较。本节对 CMS 13 TeV 双轻子道 $t\bar t$ 微分截面测量~\cite{Sirunyan:2018ucr} 给出类似比较。这些数据是在 CT18(Z) 数据集最终冻结之后发布的。QCD NNLO 理论预言用 \texttt{fastNNLO} 表~\cite{fastnnlo:grids} 配 CT18 NNLO PDF 获得。此处获取理论预言所用顶夸克质量为 $m_t^{\rm pole}=172.5$ GeV。
我们还展示用 CT18Z、CT18X 与 CT18A NNLO PDF 所得的理论预言，
截面 PDF 不确定度按 68\% C.L. 显示。
分布图形与数据-理论比较见图s.~\ref{pTt}-\ref{Mtt}。
在数据-理论图中，所有理论预言均按 CT18NNLO 理论归一。
误差棒代表统计误差与总系统误差的平方和。
我们观察到 $d\sigma / dp^t_T$ 与 $d\sigma / dm_{t\bar{t}}$ 二者在理论与未平移实验数据之间存在明显的斜率差。
这些差异可由数据的系统误差移动吸收：计入全部不确定度后 $\chi^2$ 良好。
我们注意到，就 $p_T$ 谱而言，用 CT18Z NNLO 得到的理论预言在大 $p_T$ 处对数据描述略好。



\begin{figure}[tb]
\begin{center}
\fbox{\parbox[c][3.2cm][c]{0.86\linewidth}{\centering [figure unavailable -- see original paper]\\ \texttt{CMS13-y-top-CT18NNLO-mt172p5\_ect.pdf}}}
\fbox{\parbox[c][3.2cm][c]{0.86\linewidth}{\centering [figure unavailable -- see original paper]\\ \texttt{CMS13-y-top-nnlo-ct18-zaxnnlo-mt172p5\_ect.pdf}}}
\caption{左：顶夸克快度 $y_t$ 分布。
右：含 CT18、CT18Z、CT18X、CT18A NNLO 的未平移数据-理论图。数据与理论预言按 CT18NNLO 理论归一。
误差棒表示统计误差与总系统误差的平方和。}
\label{yt}
\end{center}
\end{figure}



\begin{figure}[p]
\begin{center}
\fbox{\parbox[c][3.2cm][c]{0.86\linewidth}{\centering [figure unavailable -- see original paper]\\ \texttt{CMS13-mtt-CT18NNLO-mt172p5\_ect.pdf}}}
\fbox{\parbox[c][3.2cm][c]{0.86\linewidth}{\centering [figure unavailable -- see original paper]\\ \texttt{CMS13-mtt-top-nnlo-ct18-zaxnnlo-mt172p5\_ect.pdf}}}
\caption{左：顶夸克对不变质量分布。
右：含 CT18、CT18Z、CT18X、CT18A NNLO 的未平移数据-理论图。数据与理论预言按 CT18NNLO 理论归一。
误差棒为统计误差与总系统误差的平方和。}
\label{Mtt}
\end{center}
\end{figure}

电弱修正对 CT18 理论的影响示于图~\ref{MttEW}。这些修正按文献~\cite{Czakon:2017wor} 的相乘方案以 $K$ 因子形式计入，可在~\cite{EW:kfac} 获取。大的 EW 效应出现在高 $p_{T}^t$ 尾部。
不过在图中所示的 $p_T$ 范围 1--500 GeV 内，
多数情形 EW 修正不超过 3-4%。
若考虑更高 $p_T$ 区域，那里的 $K$ 因子会大得多。
EW 修正在极小程度上改善了顶夸克 $p_T$ 与 $m_{t\bar{t}}$ 分布的理论-数据符合。
用 CT18 PDF 的 NNLO QCD + NLO EW 预言之 $\chi^2/N_{pt}$ 与文献~\cite{Sirunyan:2018ucr} 表 49 所给数值相符良好。
对其余所有分布，EW 修正在所研究运动学范围内可忽略。
CT18 全球分析目前只包含 ATLAS 与 CMS 8 TeV 的 $\bar{t}t$ 微分截面测量。拟合中这些分布的 CT18 理论预言不含 EW 修正。若在拟合中纳入 EW 修正，鉴于当前考虑分布的运动学范围内 EW 修正的大小，其对被拟合 PDF 的影响可以忽略。


\begin{figure}
\begin{center}
\fbox{\parbox[c][3.2cm][c]{0.86\linewidth}{\centering [figure unavailable -- see original paper]\\ \texttt{CMS13-pt-top-qcdnnloewnlo-ct18nnlo-mt172p5\_ect.pdf}}}
\fbox{\parbox[c][3.2cm][c]{0.86\linewidth}{\centering [figure unavailable -- see original paper]\\ \texttt{CMS13-mtt-qcdnnloewnlo-ct18nnlo-mt172p5\_ect.pdf}}}
\caption{NLO EW 修正的影响。顶夸克 $p_T$ 分布与 $t\bar{t}$ 对不变质量分布的未平移数据-理论图。
误差棒为统计误差与总系统误差的平方和。红色带虚线误差棒的数据点表示数据除以含 NLO EW 修正的 CT18NNLO 理论。黑色带实线误差棒的数据点表示按 CT18NNLO 理论归一的数据。CT18ZNNLO 理论预言（黑色虚线误差带）也按 CT18NNLO 归一。为便于查看，数据 vs (理论+EW) 点在同一分箱内向右稍作移动。}
\label{MttEW}
\end{center}
\end{figure}


\begin{figure}[h]
	\fbox{\parbox[c][3.2cm][c]{0.86\linewidth}{\centering [figure unavailable -- see original paper]\\ \texttt{pdfs\_CT18NNLO-58\_CT18pCMS13ytt\_100-0GeV\_S90CL\_\_00\_\_\_glu\_\_pdfr\_cus-lin.pdf}}}
	\fbox{\parbox[c][3.2cm][c]{0.86\linewidth}{\centering [figure unavailable -- see original paper]\\ \texttt{pdfs\_CT18ZNNLO-58\_CT18ZpCMS13ytt\_100-0GeV\_S90CL\_\_00\_\_\_glu\_\_pdfr\_cus-lin.pdf}}}
	\caption{CMS 13 TeV $t \bar t$ 数据的 $y_{t \bar t}$ 微分截面测量对 CT18（左）与 CT18Z（右）胶子 PDF 的影响。
	\label{fig:cms13ytt}}
\end{figure}

在各种一维 $t \bar{t}$ 微分分布中，顶夸克对快度分布 $y_{t \bar t}$ 显示 CMS 数据与 CT18 预言符合良好。为考察该数据会如何修改 CT18(Z) 胶子 PDF，我们用 \texttt{ePump} 程序~\cite{Hou:2019gfw} 在把 CMS 13 TeV $y_{t \bar t}$ 数据纳入拟合后更新 CT18(Z) PDF。
如图~\ref{fig:cms13ytt} 所示，两种情形下更新后的胶子 PDF 误差带（记作 CT18pCMS13ytt）在 $x$ 从 0.1 到 0.4 都极轻微地收窄。
关于这些数据集的进一步讨论将另行发表。



\subsection{高 $x$ Drell-Yan 预言
\label{sec:SeaQuest}
}
固定靶 Drell-Yan 测量为核子（与核）PDF 的 $x$ 依赖性提供重要探针。这一事实推动了包括费米实验室 E866/NuSea
实验~\cite{Towell:2001nh} 在内的一系列实验；该实验测定了直至较大 $x_2$（靶部分子动量分数）的归一化氘核-质子截面比
$\sigma_{pd} \big/ 2\sigma_{pp}$。从领头阶夸克-部分子模型分析可见，该比值预期对 PDF 比 $\bar{d}/\bar{u}$ 的 $x$ 依赖性特别敏感，
使其成为研究轻夸克海味对称破缺的有利可观测量。如 Sec.~\ref{sec:moments} 讨论 Gottfried 求和规则时所述，$\mathrm{SU}(2)$ 对称破缺被认为具有非微扰起源。

\begin{figure}[h]
	\begin{center}
		\fbox{\parbox[c][3.2cm][c]{0.86\linewidth}{\centering [figure unavailable -- see original paper]\\ \texttt{CT18-prediction\_E866\_68.pdf}}}
	\end{center}
	\vspace{-2ex}
	\caption{ 基于 CT18（黑色外带）与 CT18Z（绿色内带）、针对费米实验室 SeaQuest 实验~\cite{Aidala:2017ofy} 将探测的较大 $x_2 \gtrsim 0.1$ 区域固定靶 Drell-Yan
	截面 $\sigma_{pd} \big/ 2\sigma_{pp}$ 的理论预言。为便于比较，也绘出较早 E866 数据~\cite{Towell:2001nh} 的较高 $x_2$ 部分（蓝色菱形）。 }
\label{fig:hix_DY}
\end{figure}

引人注目的是，E866~\cite{Towell:2001nh} 发现有证据表明当 $x_2$ 接近并超过 $x \gtrsim 0.25$ 时截面比跌破 1，即 $\sigma_{pd} \big/ 2\sigma_{pp} < 1$，见图~\ref{fig:hix_DY} 中 E866 比值点较高的 $x_2$ 部分。这一事实令诸多理论模型感到意外。
因此 E866 结果激发了以更高精度把类似测量推进到更大 $x_2$ 的兴趣——这正是其后费米实验室 SeaQuest/E906 实验~\cite{Aidala:2017ofy} 的主要目标，结果预计很快公布。为此，我们在图~\ref{fig:hix_DY} 中展示基于更新后的 CT18（黑带）与 CT18Z（绿带）全球分析、按 68\% C.L. 外推到超出 E866 所探范围的更高 $x_2$ 的理论预言。虽然 CT18
与 CT18Z 受 E866 比值数据约束，氘核-质子比的理论预言直到
$x\! <\! 0.4$ 仍保持在 1 之上或与之相符。高 $x$ 区更多精密数据将对厘清截面比行为及其对核子海的蕴含起到关键作用。

\clearpage
